\subsection{Application: change of variable in the Lebesgue integral}

Let I be an interval (bounded or not) of R, a its left end-point and b its right end-point in \(\overline{R}\) , and \(\mu\) the Lebesgue measure on I. For every \(\mu\) -integrable function f and every interval \(H \subset I\) , with left end-point \(\alpha\) and right end-point \(\beta\) , we write \(\int_{\alpha}^{\beta} \mathbf{f}(t) dt\) instead of \(\int_{H} \mathbf{f}(t) dt = \int_{H} \mathbf{f} d\mu\) , and we set \(\int_{\beta}^{\alpha} \mathbf{f}(t) dt = -\int_{\alpha}^{\beta} \mathbf{f}(t) dt\) ; the meaning thus given to these symbols coincides with that attributed to them in FRV, II, §§1 and 2, when f is a regulated function with compact support (Ch. IV, §4, No.~4, Example).

Let \(g\) be a numerical function defined on \(I\) and locally \(\mu\) -integrable, \(x_0\) a point of \(I\) ; for every \(x \in I\) , set

\[
G (x) = c + \int_ {x _ {0}} ^ {x} g (t) d t \quad (c \text {   a   constant }).\tag{11}
\]

The numerical function G is continuous on I; this follows at once from Lebesgue's theorem (Ch. IV, §4, No.~3, Cor. 1 of Th. 2), since the product of g and the characteristic function of the interval with end-points x and \(x + h\) tends to a negligible function as h tends to 0. Therefore G(I) is an interval of R. Throughout this No., we will regard G as a mapping of I onto the locally compact space G(I). We denote by \(\lambda\) the measure \(g \cdot \mu\) on I.

Suppose first that g is \(\mu\) -integrable. Then, the same reasoning as above shows that the limits \(\mathrm{G}(a+)\) and \(\mathrm{G}(b-)\) exist and are finite; moreover, the measure \(|\lambda|\) is bounded ( \(\S5\) , No.~3, Cor. of Th. 1), and the mapping G of I into \(\mathrm{G}(\mathrm{I})\) is \(\lambda\) -proper.

PROPOSITION 8. --- Suppose g is \(\mu\) -integrable. If J denotes the open interval in R with end-points \(\mathrm{G}(a+)\) and \(\mathrm{G}(b-)\) , the image under G of the measure \(g \cdot \mu\) is the measure \(\varphi_{J} \cdot \nu\) if \(\mathrm{G}(a+) \leqslant \mathrm{G}(b-)\) and the measure \(-\varphi_{J} \cdot \nu\) if \(\mathrm{G}(a+) \geqslant \mathrm{G}(b-)\) (where \(\nu\) denotes Lebesgue measure on \(\mathrm{G}(\mathrm{I})\) ).

It suffices to prove that, for every function \(f \in \mathcal{K}\big(\mathrm{G}(\mathrm{I})\big)\) ,

\[
\int_ {\mathrm{G} (a +)} ^ {\mathrm{G} (b -)} f (\xi) d \xi = \int_ {a} ^ {b} f (\mathrm{G} (t)) g (t) d t.\tag{12}
\]

Now, this formula has already been proved for \(g \in \mathcal{K}(\mathrm{I})\) (FRV, II, §2, No.~1, formula (1)). Let us pass to the general case; there exists a sequence \((g_n)\) of functions in \(\mathcal{K}(\mathrm{I})\) such that: \(1^\circ\) the sequence \((g_n(t))\) tends to \(g(t)\) almost everywhere in I; \(2^\circ\) there exists a \(\mu\) -integrable function \(h \geqslant 0\) such that \(|g_n| \leqslant h\) for all \(n\) (Ch. IV, §3, No.~4, Th. 3). It follows at once from Lebesgue's theorem that, setting \(G_n(x) = c + \int_{x_0}^x g_n(t) dt\) , the sequence \((G_n)\) converges uniformly to G on I, and that the numbers \(G_n(a+)\) and \(G_n(b-)\) tend respectively to \(G(a+)\) and \(G(b-)\) . Let \(f'\) be a function in \(\mathcal{K}(\mathbf{R})\) that extends \(f\) ; the foregoing proves that \(f'(G_n(t))\) tends to \(f'(G(t)) = f(G(t))\) for all \(t \in \mathrm{I}\) ; applying Lebesgue's theorem, one sees that the formula (12) results from the formula

\[
\int_ {\mathrm{G} _ {n} (a +)} ^ {\mathrm{G} _ {n} (b -)} f ^ {\prime} (\xi) d \xi = \int_ {a} ^ {b} f ^ {\prime} \bigl (\mathrm{G} _ {n} (t) \bigr) g _ {n} (t) d t
\]

by passage to the limit.

COROLLARY. --- If a function f defined on G(I), with values in \(\overline{R}\) or in a Banach space, is such that the function \(t \mapsto \mathbf{f}(G(t))g(t)\) is integrable on I for Lebesgue measure, then f is integrable on J for Lebesgue measure and

\[
\int_ {\mathrm{G} (a +)} ^ {\mathrm{G} (b -)} \mathbf {f} (\xi) d \xi = \int_ {a} ^ {b} \mathbf {f} (\mathrm{G} (t)) g (t) d t\tag{13}
\]

(formula for change of variable in the Lebesgue integral).

For, \(\mathbf{f}\big(\mathrm{G}(t)\big)\) is integrable for the measure \(|g|\cdot\mu\) , hence also for the measures \(g^{+}\cdot\mu\) and \(g^{-}\cdot\mu\) ; it follows from Th. 1 (No.~2) that f is integrable for the image measures \(\mathrm{G}(g^{+}\cdot\mu)\) and \(\mathrm{G}(g^{-}\cdot\mu)\) , hence also for the measure \(\varphi_{J}\cdot\nu\) , and that (13) holds, on taking into account Prop. 8 and formula (9).

It can happen that \(\mathbf{f}\) is integrable on J for Lebesgue measure, but that \(t \mapsto \mathbf{f}(\mathbf{G}(t))g(t)\) is not integrable on I for Lebesgue measure (Exer. 10).

Now suppose that g maintains a constant sign almost everywhere (and is locally \(\mu\) -integrable); one may suppose for example that \(g(t) \geqslant 0\) almost everywhere in I. Then G is an increasing continuous function on I, therefore \(\mathrm{G}(a+)\) and \(\mathrm{G}(b-)\) exist (but may be infinite). Moreover, G is a \(\lambda\) -proper mapping of I into \(\mathrm{G}(\mathrm{I})\) : for, if \(\mathrm{G}(b-) \in \mathrm{G}(\mathrm{I})\) , there is an \(x_{1} \geqslant x_{0}\) such that G is constant for \(x \geqslant x_{1}\) , and then the inverse image under G of the compact interval \([G(x_{0}), G(b-)]\) is \(\lambda\) -integrable; if, on the contrary, \(G(b-) \notin G(I)\) , then the inverse image under G of every compact interval with left end-point \(G(x_{0})\) , contained in G(I), differs from a compact interval by at most a \(\lambda\) -negligible interval. One argues similarly for the compact intervals with right end-point \(G(x_{0})\) , whence our assertion. Moreover:

PROPOSITION 9. --- Suppose \(g \geqslant 0\) and locally \(\mu\) -integrable. Then, the image under \(G\) of the positive measure \(g \cdot \mu\) is Lebesgue measure on \(G(I)\) . For a function \(f\) , defined on \(G(I)\) , with values in \(\overline{R}\) or in a Banach space, to be integrable on \(G(I)\) for Lebesgue measure, it is necessary and sufficient that the function \(t \mapsto f(G(t))g(t)\) be integrable on \(I\) for Lebesgue measure, in which case the relation (13) holds.

The first part of the statement follows from the fact that the formula (12) is valid for every function \(f \in \mathcal{K}(\mathrm{G}(\mathrm{I}))\) ; for, the support of the function \(t \mapsto f(\mathrm{G}(t))\) is contained in an interval \(K \subset I\) on which g is integrable, by virtue of the above remarks, and it suffices to apply Prop. 8 to K. The second part is a consequence of Th. 1 of No.~2.

